% Lecture notes accompanying the "Repair or Replace" visualization (SML, ETH Zurich).
\documentclass[11pt,a4paper]{article}

\usepackage[margin=2.7cm]{geometry}
\usepackage[T1]{fontenc}
\usepackage{amsmath,amssymb,amsthm,mathtools}
\usepackage{lmodern}
\usepackage{microtype}
\usepackage{booktabs}
\usepackage{xcolor}
\definecolor{inknavy}{RGB}{18,52,98}
\usepackage[colorlinks=true,linkcolor=inknavy,urlcolor=inknavy,citecolor=inknavy]{hyperref}
\hypersetup{pdftitle={Reinforcement Learning: Lecture Notes},pdfauthor={Carlos Cotrini}}
\urlstyle{same}

% Definitions, examples, and remarks share one counter, numbered within sections.
\theoremstyle{definition}
\newtheorem{definition}{Definition}[section]
\newtheorem{example}[definition]{Example}
\theoremstyle{remark}
\newtheorem{remark}[definition]{Remark}
\numberwithin{equation}{section}

\newcommand{\E}{\mathbb{E}}
\newcommand{\sa}[1]{\textsc{#1}}   % state and action names, e.g. \sa{healthy}
\newcommand{\vizurl}{\url{https://sml-fs26.github.io/classic_rl/repair-or-replace/}}
\newcommand{\vizhref}[1]{\href{https://sml-fs26.github.io/classic_rl/repair-or-replace/}{#1}}
\DeclareMathOperator*{\argmax}{arg\,max}

\title{Reinforcement Learning\\[0.5ex]\large Lecture Notes}
\author{Carlos Cotrini\\ETH Zürich}
\date{\today}

\begin{document}
\maketitle

\noindent
These notes accompany the interactive visualization \emph{Repair or Replace}, available at
\begin{center}
\vizurl
\end{center}
\noindent
The concepts developed below map one-to-one onto its scenes.

\medskip

% 1: the hook -- sequential decisions under uncertainty.
\section{Motivation}

Many decisions a manager faces are sequential. What is decided today changes what must be
decided tomorrow, and the consequences unfold over weeks and months rather than at once. The
outcomes are uncertain on top: the same decision may turn out well one time and badly the
next. Picture the scenario behind these notes. You operate a delivery business with a single
van, and each week you make exactly one call about that machine: keep running it and collect
the week's profit, at the risk of a breakdown and of gradual wear; repair it, paying for a
week in the shop so that it fails less often afterwards; or replace it outright, a large
expense that buys a fresh start. Every option trades a short-term cost against a long-term
payoff, and chance interferes at every step.

Reinforcement learning is a framework for computing decision strategies that maximize the
reward accumulated over time in precisely such problems. One running example, the
repair-or-replace problem of the visualization, carries everything.

% 2: the model.
\section{Markov decision processes}

The first step is a precise model of the decision problem. The standard model is the
\emph{Markov decision process}.\footnote{The name records that the next state depends only on
the current state and action, not on the past.}

\begin{definition}[Markov decision process]
A \emph{Markov decision process} (MDP) is a quadruple, a package of four ingredients,
$(S, A, T, \gamma)$, where $S$ is the set of \emph{states} the system can be in; $A$ is the
set of \emph{actions} available to the decision maker; $T$ is the \emph{transition dynamics},
which describes how the system reacts when an action is taken; and $\gamma \in [0,1]$, the
\emph{discount factor}, is a number between $0$ and $1$. ($\gamma$ is the Greek letter
gamma.)
\end{definition}

The sets $S$ and $A$ are defined a priori: they are fixed in advance, part of the problem
description, and writing them down is the modeler's first task. The transition dynamics is
then defined on top of them. It is a \emph{randomized function} $T$ that consumes a state $s$
and an action $a$ and produces a new state $s'$ and a reward $r$:
\[
(s', r) = T(s, a).
\]
In words: feed the current state and the chosen action into the dynamics, and out come the
state the system lands in next and the reward, a number measuring the immediate gain or loss
of that step.

\begin{remark}\label{rem:randomdyn}
The transition dynamics is \emph{random}: feeding the same state and action $(s, a)$ in twice
may produce different new states and different rewards.\footnote{Scene ``You run the fleet''
of the interactive \vizhref{visualization} lets you sample these dynamics.}
\end{remark}

The discount factor $\gamma$ is just a number between $0$ and $1$; its use will be explained
later.

\begin{example}[the repair-or-replace MDP]
You manage a one-van delivery fleet; the visualization calls the van \emph{Old Bessie}. Each
week the van sits at one of four wear levels, and you make exactly one call:
\[
S = \{\sa{healthy}, \sa{worn}, \sa{shaky}, \sa{failing}\},
\qquad
A = \{\sa{run}, \sa{service}, \sa{replace}\}.
\]
In words: four wear levels, from factory-fresh to nearly dead, and three possible calls
(\sa{service} is the ``repair'' of the title). Rewards are weekly net money, in plain
units. The numbers governing \sa{run} are:
\begin{center}
\begin{tabular}{lcc}
\toprule
state & weekly profit & breakdown odds \\
\midrule
\sa{healthy} & $+95$ & $2\%$  \\
\sa{worn}    & $+72$ & $8\%$  \\
\sa{shaky}   & $+40$ & $28\%$ \\
\sa{failing} & $+16$ & $55\%$ \\
\bottomrule
\end{tabular}
\end{center}
Under \sa{run} the van drives the route and earns the profit in the table, unless it breaks
down, with the odds in the table: the week then books a single net reward equal to the profit
minus a breakdown cost of $280$ (from \sa{healthy}: $95 - 280 = -185$), and the van ends the
week at \sa{failing}. Even without a breakdown, wear tends to slip by one or two levels over
the weeks. Under \sa{service} the van spends the week in the shop: reward $-50$, no
deliveries, and the wear often improves; from \sa{worn}, the van returns to \sa{healthy}
$95\%$ of the time. Under \sa{replace} the reward is $-130$, and the next week starts
with a
brand-new van at \sa{healthy}, whatever the old one looked like. The dynamics is random in
the sense of Remark~\ref{rem:randomdyn}: feeding $(\sa{worn}, \sa{run})$ into $T$ may return
$(\sa{shaky}, +72)$ one week and $(\sa{failing}, -208)$ another. There is no final week (the
fleet runs indefinitely), and the visualization fixes $\gamma = 0.9$.
\end{example}

% 3: strategies and the histories they generate.
\section{Policies and trajectories}

An MDP describes the world the manager operates in; it does not say how to operate. That is
the job of a policy.

\begin{definition}[policy]
A \emph{policy} defines how to act at each state. Formally, a policy is a function
$\pi \colon S \to A$: it takes a state $s$ and returns the action $\pi(s)$ to be played
whenever the system is in $s$. ($\pi$ is the Greek letter pi.)
\end{definition}

Executing a policy on the MDP, week after week, generates a record of everything that
happened. That record is the central object of these notes.

\begin{definition}[trajectory]
The \emph{trajectory} of a policy $\pi$ is the sequence
\[
s_1,\; a_1,\; r_1,\; s_2,\; a_2,\; r_2,\; \dots,\; s_i,\; a_i,\; r_i,\; \dots
\]
produced by executing $\pi$ on the MDP starting from state $s_1$. In words: states, actions,
and rewards alternate, and the subscript counts the weeks. The mechanics are as follows. The
first action results from applying the policy to the first state, $a_1 = \pi(s_1)$. The first
reward and the second state result from applying the transition dynamics to that pair,
$(s_2, r_1) = T(s_1, a_1)$. Then $a_2 = \pi(s_2)$, the dynamics produce
$(s_3, r_2) = T(s_2, a_2)$, and so on without end.
\end{definition}

\begin{remark}\label{rem:randomtraj}
The trajectory is \emph{random}: starting again from the same $s_1$ and executing the same
policy may produce a different trajectory, because the transition dynamics is random (and so,
possibly, is the policy, if one allows policies that randomize their choices).
\end{remark}

\begin{example}[a trajectory of the fleet]\label{ex:trajectory}
Consider the playbook $\pi$ with $\pi(\sa{healthy}) = \sa{run}$,
$\pi(\sa{worn}) = \sa{service}$, and $\pi(\sa{shaky}) = \pi(\sa{failing}) = \sa{replace}$.
One execution of $\pi$ from $s_1 = \sa{healthy}$ produced\footnote{Scene ``The trajectory''
of the interactive \vizhref{visualization} replays such a tape week by week.}
\begin{align*}
s_1 &= \sa{healthy}, & a_1 &= \sa{run},     & r_1 &= +95,\\
s_2 &= \sa{worn},    & a_2 &= \sa{service}, & r_2 &= -50,\\
s_3 &= \sa{healthy}, & a_3 &= \sa{run},     & r_3 &= +95,\\
s_4 &= \sa{healthy}, & a_4 &= \sa{run},     & r_4 &= +95,
\end{align*}
after which $s_5 = \sa{worn}$ and the walk continues, one week per step, forever. In words:
the van started healthy and earned $95$ but slipped to worn; a $50$ service restored it; two
profitable runs followed; by week five it was worn again. Remark~\ref{rem:randomtraj} is at
work here: another execution of the same $\pi$ from the same $s_1$ might open with a
first-week breakdown, $s_1 = \sa{healthy}$, $a_1 = \sa{run}$, $r_1 = -185$,
$s_2 = \sa{failing}$.
\end{example}

% 4: scoring a trajectory, and what to maximize.
\section{The goal: accumulating reward}

A policy is good if it produces trajectories that accumulate as much reward as possible. The
first task is to make ``accumulate'' precise.

\begin{definition}[cumulative reward]
Take one particular trajectory $s_1, a_1, r_1, s_2, a_2, r_2, \dots$. Its \emph{cumulative
reward}, also called its \emph{return}, is
\begin{equation}\label{eq:return}
G_1 \;=\; r_1 + \gamma\, r_2 + \gamma^2 r_3 + \cdots + \gamma^{\,i-1} r_i + \cdots.
\end{equation}
In words: add up the rewards of all weeks, forever (that is what the dots mean), scaling the
reward of week $i$ by $\gamma^{\,i-1}$: week one counts in full, week two is scaled by
$\gamma$, week three by $\gamma^2$, and so on. ($G$ stands for \emph{gain}.)
\end{definition}

Here the discount factor plays its very important role: it specifies how much we care about
future rewards. With $\gamma$ close to $0$, the scaling factors die out almost immediately
and only near-term rewards matter; with $\gamma$ close to $1$, they fade slowly and rewards
far in the future count almost as much as immediate ones. Managers will recognize the device:
it is the discounting of future cash flows familiar from net-present-value calculations, and,
as with the NPV of a perpetuity, the shrinking factors keep the never-ending total finite.
With the visualization's $\gamma = 0.9$, next week's reward counts $90\%$ of this week's, the
week after $81\%$, and so on.\footnote{Scene ``Return over the van's life'' of the
interactive \vizhref{visualization} has a slider for $\gamma$.} For the trajectory of
Example~\ref{ex:trajectory}, formula~\eqref{eq:return} gives
$G_1 = 95 + 0.9 \cdot (-50) + 0.9^2 \cdot 95 + 0.9^3 \cdot 95 + \cdots
= 95 - 45 + 76.95 + 69.255 + \cdots$.

The same accounting can be started at any week, not only the first. For an arbitrary week
$i$, define $G_i$ as the cumulative reward from week $i$ on:
\begin{equation}\label{eq:gi}
G_i \;=\; r_i + \gamma\, r_{i+1} + \gamma^2 r_{i+2} + \cdots.
\end{equation}
In words: $G_i$ is what the trajectory accumulates from week $i$ onward, with the discounting
restarted there. Comparing \eqref{eq:gi} at weeks $i$ and $i+1$ yields a recurrence that
carries the rest of these notes:
\begin{equation}\label{eq:recurrence}
G_i \;=\; r_i + \gamma\, G_{i+1}.
\end{equation}
In words: what is accumulated from week $i$ on equals this week's reward plus $\gamma$ times
what is accumulated from week $i+1$ on.

\begin{remark}
$G_i$ is \emph{random}. Because of the randomness of the transition dynamics, $G_i$ comes out
differently after every execution of the policy starting from state $s_i$.
\end{remark}

A quantity that changes from run to run cannot be maximized directly; we average it instead.

\begin{definition}[expectation]\label{def:expectation}
The \emph{expectation} $\E[G_i]$, read ``the expected value of $G_i$,'' is the average of
$G_i$ over a sufficiently large set of trajectories: execute the policy many times, always
starting from the same first state $s_1$ (for the fleet: a \sa{healthy} van), record the
value of $G_i$ of each execution, and average the records.
\end{definition}

We can now state the ultimate objective of reinforcement learning: \emph{find the policy
$\pi$ that maximizes $\E[G_1]$}, the policy whose trajectories accumulate, on average, the
most reward. A policy achieving this maximum is called \emph{optimal}. Note the contrast
with $G_i$ itself: $\E[G_i]$ is a single real number that does not depend on the
trajectory, because the averaging has removed the randomness. Therefore
$\E[G_i]$ is \emph{not} random.

% 5: the scorecard and the equation that pins it down.
\section{The \texorpdfstring{$Q$}{Q}-function and the Bellman equation}\label{sec:qfunction}

How can the best policy be discovered? The classical route goes through an auxiliary function
that scores every pair of a state and an action.

\begin{definition}[the $Q$-function]
For a state $s$ and an action $a$, $Q^\star(s, a)$ is the cumulative reward one can expect
when executing an optimal policy from state $s$, where the \emph{first} action taken there is
$a$. In notation,
\[
Q^\star(s, a) \;=\; \E\bigl[\, G_i \bigm| s_i = s,\ a_i = a \,\bigr].
\]
In words: the expected value of $G_i$ conditioned on $s_i$ being $s$ and $a_i$ being $a$ (the
vertical bar is read ``given that''); that is, the average of $G_i$ over many executions,
counting only those executions that are in state $s$ at week $i$ and take action $a$ there as
a one-off, even when an optimal policy would choose differently, acting optimally afterwards.
(The star in $Q^\star$ marks optimality.)
\end{definition}

Which week $i$? It does not matter: the fleet has no final week, so it looks the same from
week $3$ as from week $30$, and every week gives the same average.

\begin{remark}
$Q^\star(s, a)$ is different from $\E[G_i]$. The number $\E[G_i]$ averages over \emph{all}
executions of the policy, whatever state they happen to be in at week $i$. $Q^\star(s, a)$
restricts the average to the executions that pass through $s$ at week $i$ and are made to
play $a$ there, acting optimally afterwards. And while $\E[G_i]$ is a single number,
$Q^\star$ is a whole table of numbers, one per state-action pair: for the fleet,
$4 \times 3 = 12$ entries.
\end{remark}

The table is worth the trouble because the optimal policy can be read off it: at state $s$,
the optimal action is simply the action $a$ that maximizes $Q^\star(s, a)$. In notation,
\[
\pi^\star(s) \;=\; \argmax_{a}\, Q^\star(s, a).
\]
In words: in state $s$, compare the values $Q^\star(s, a)$ of all available actions $a$ and
play an action with the largest value; $\argmax$ denotes exactly that action, and
$\pi^\star$ denotes the optimal policy.

It remains to compute the table. Suppose an optimal policy produced the trajectory
$\dots, s_i, a_i, r_i, s_{i+1}, a_{i+1}, \dots$. Averaging both sides of the
recurrence~\eqref{eq:recurrence}, $G_i = r_i + \gamma\, G_{i+1}$, over many such executions,
it must be the case that the \emph{Bellman equation} holds:
\begin{equation}\label{eq:bellman}
\boxed{\;Q^\star(s_i, a_i) \;=\; \E\bigl[\, r_i + \gamma\, Q^\star(s_{i+1}, a_{i+1}) \,\bigr]\;}\,.
\end{equation}
In words: the value of this week's state and action equals this week's reward plus $\gamma$
times the value of next week's state and action, averaged over the randomness of the week.
Replacing $G_{i+1}$ by its average $Q^\star(s_{i+1}, a_{i+1})$ is legitimate: averaging in
two rounds (first among the executions that share next week's state and action, then across
the week's outcomes) gives the same result, just as a company-wide average can be computed
branch by branch.

We can therefore find $Q^\star$ by solving the Bellman equations, one equation per
state-action pair, twelve for the fleet. Since the optimal policy at $s_{i+1}$ picks the
action with the largest $Q^\star$-value, the right-hand side of \eqref{eq:bellman} involves
the same unknowns $Q^\star(s, a)$ as the left-hand side, so this is a system of equations in
the entries of the table. The visualization solves it iteratively; its $Q^\star$-scorecard
and Bellman-check scenes display the solution for the fleet at
$\gamma = 0.9$.\footnote{Scenes ``$Q^\star$: the action scorecard'' and ``The Bellman
equation,'' interactive in the \vizhref{visualization}.} Two entries tell the story:
$Q^\star(\sa{shaky}, \sa{service}) = 126.4$, while
$Q^\star(\sa{shaky}, \sa{replace}) = 149.9$. The optimal playbook replaces a van that still
starts every morning.

% 6: when nobody prints the odds, learn the table from the lived weeks.
\section{SARSA: learning from the logbook}\label{sec:sarsa}

Solving the Bellman equations requires knowing the odds. The average
in~\eqref{eq:bellman} runs over the randomness of the week, and computing an average over
outcomes means knowing how likely each outcome is: the transition dynamics $T$ itself. In
the visualization those odds are printed on the fleet sheet; in reality nobody hands them
to you, because no manual lists the breakdown odds for \emph{your} van on \emph{your}
routes. And for realistic fleets the table of states explodes anyway (count odometer
readings, age, weeks since the last service, route load, then forty vans, and the
visualization's illustrative tally reaches some four million states for a single van and
half a billion scorecard cells for a modest fleet).\footnote{Scene ``Why DP doesn't
scale'' of the interactive \vizhref{visualization} prices this out.} What the manager does
have is the \emph{logbook}, week after week of lived experience: the state the van was
in, the call made, the money booked, the state it landed in, and the call made there.

SARSA needs nothing else. Maintain a table $q(s, a)$ of \emph{estimates}, current best
guesses of $Q^\star(s, a)$, one cell per state-action pair, twelve for the fleet, with
every cell starting at $0$, as in the visualization. Then, after every lived week, improve
the single cell just used.

The evidence for that improvement comes from the recurrence~\eqref{eq:recurrence}: what is
accumulated from week $i$ on equals this week's reward plus $\gamma$ times what is
accumulated from week $i+1$ on. The future itself has not happened yet, and waiting for
it would take forever, so the table's current best guess for the next pair stands in for
it: the same move as plugging a forecasted cash flow into an NPV calculation. One week of
experience $(s_i,\, a_i,\, r_i,\, s_{i+1},\, a_{i+1})$, one logbook line, therefore reports: ``this
week, the value of the pair $(s_i, a_i)$ looked like
$r_i + \gamma\, q(s_{i+1}, a_{i+1})$,'' that is, this week's money plus $\gamma$ times the
current estimate of the pair the van moved on to. SARSA folds that report into the one
cell it concerns:
\begin{equation}\label{eq:sarsa}
q(s_i, a_i) \;\leftarrow\; q(s_i, a_i)
\;+\; \alpha \bigl[\, r_i + \gamma\, q(s_{i+1}, a_{i+1}) - q(s_i, a_i) \,\bigr].
\end{equation}
In words: the arrow is an instruction, not an equation; compute the right side and replace
the left side by it. Three steps unpack the instruction. First, name the bracket: it is
the \emph{surprise}, the gap between what this week's evidence says the pair is worth and
what the table currently believes. A positive surprise means the week looked better than
the table believed; a negative one, worse. Second, the move: shift the belief a small
fraction $\alpha$ of the way toward the evidence. The number $\alpha$ is the
\emph{learning rate}, $0.15$ in the visualization. ($\alpha$ is the Greek letter alpha.)
Never overwrite, only nudge: one lucky week must not erase a year of evidence. To see one
nudge land, feed in week one of Example~\ref{ex:trajectory} with the table still all
zeros: the logbook line is $(\sa{healthy}, \sa{run}, +95, \sa{worn}, \sa{service})$; the
evidence is $95 + 0.9 \cdot q(\sa{worn}, \sa{service}) = 95 + 0.9 \cdot 0 = 95$; the
surprise is $95 - 0 = 95$; and the cell $q(\sa{healthy}, \sa{run})$ moves from $0$ to
$0 + 0.15 \cdot 95 = 14.25$. One cell moves; the other eleven stay put. Third, the
punchline. The Bellman
equation~\eqref{eq:bellman} says that $Q^\star(s_i, a_i)$ is the \emph{average} of the
evidence $r_i + \gamma\, Q^\star(s_{i+1}, a_{i+1})$ over many weeks, and an average over
many executions is exactly how the expectation was defined in
Definition~\ref{def:expectation}. SARSA computes that average the only way available
without the odds: incrementally, as the weeks arrive. The update rule is the standard
recipe for a running average (Appendix~\ref{app:runavg} derives it). Set \eqref{eq:sarsa}
beside~\eqref{eq:bellman}: the update rule is the Bellman equation with the expectation
replaced by an arriving average and $Q^\star$ replaced by the current guess $q$. SARSA
solves the Bellman equation by averaging evidence instead of computing expectations.

\begin{remark}
The name SARSA spells the five logbook entries each update consumes: State, Action,
Reward, State, Action.
\end{remark}

Which calls should be made while the table is still being learned? Mostly the one the
table currently favors: at state $s$, play the action with the largest $q(s, a)$. But in a
small random share of the weeks (a fraction $\varepsilon = 0.15$ in the visualization,
held constant, like $\alpha$; $\varepsilon$ is the Greek letter epsilon) the call is
instead drawn at random, because a cell that is never visited is never improved: you
cannot learn the value of servicing if you never service. The $a_{i+1}$ in the update is
the action \emph{actually} chosen at $s_{i+1}$ by
this very playbook, so SARSA learns the value of what it really does, occasional
experiments included.

Run this, week after week, on one endless stream of lived weeks: drive, get billed, nudge
one cell. The table drifts toward $Q^\star$, and reading off the best action in each row
recovers the playbook. In the visualization the same three bands as in the scorecard of
Section~\ref{sec:qfunction} (\sa{run} when \sa{healthy}, \sa{service} when \sa{worn},
\sa{replace} when \sa{shaky} or \sa{failing}) emerge within a few hundred weeks of raw
experience, the odds never shown.\footnote{Scene ``SARSA: learn from the logbook'' of the
interactive \vizhref{visualization}: watch one week nudge one cell.} One
honest caveat: with $\alpha$ and $\varepsilon$ held constant, the twelve numbers
themselves stay rough; the experiments cost money, so the cells settle below the
$Q^\star$ values and keep jittering. But which call tops each row, and that is all the
playbook reads, comes out right.

% Closing: recap and invitation.
\section*{Closing}

These notes modeled sequential decision problems as Markov decision processes, let policies
generate random trajectories, scored each trajectory by its cumulative reward, averaged that
score to obtain something maximizable, and characterized the best policy through the
$Q$-function and the Bellman equation. When the odds are unknown, SARSA learns the same
scorecard from the logbook alone. The fastest way to make all of this concrete is to
play: open \vizurl, run the fleet for a quarter, drag the $\gamma$ slider, and watch the
$Q^\star$ scorecard fill in, scene by scene, in the same order as these notes.

% Appendices: the averaging machinery behind the update rule.
\appendix

\section{The Robbins--Monro algorithm}\label{app:rm}

Strip the fleet away and a bare problem remains: find the number $x^\star$ that
\emph{equals the average} of noisy evidence, when the average cannot be computed directly
and the evidence arrives one sample at a time. The Robbins--Monro algorithm keeps a running
guess $x$ and, when the $n$-th sample arrives, nudges it:
\[
x \;\leftarrow\; x + \alpha_n\,(\text{sample} - x).
\]
In words: move the guess a fraction $\alpha_n$ of the way toward each sample as it
arrives. The step sizes $\alpha_1, \alpha_2, \alpha_3, \dots$ must satisfy two demands.
The nudges must add up without bound, meaning that
$\alpha_1 + \alpha_2 + \alpha_3 + \cdots$ grows past any number, so that the iteration can
travel from any bad start all the way to $x^\star$. And their squares must add up to
something finite, meaning that $\alpha_1^2 + \alpha_2^2 + \alpha_3^2 + \cdots$ stays
finite, so that the noise in the samples eventually cancels out instead of jiggling the
guess forever. The choice $\alpha_n = 1/n$ satisfies both demands: the running total
$1 + 1/2 + 1/3 + \cdots$ climbs past any bound, just ever more slowly (a classical, if
surprising, fact), while the squared total $1 + 1/4 + 1/9 + \cdots$ settles below $2$.

The SARSA update~\eqref{eq:sarsa} is exactly a Robbins--Monro iteration applied cell by
cell to the Bellman equation~\eqref{eq:bellman}: the guess is the cell $q(s_i, a_i)$, and
the arriving sample is the evidence $r_i + \gamma\, q(s_{i+1}, a_{i+1})$. With one honest
extra: that sample is itself measured with the current table $q$, so the algorithm pulls
itself up by its own bootstraps, and making this rigorous took decades. And a second
honest extra: Section~\ref{sec:sarsa} holds $\alpha$ constant at $0.15$, deliberately
failing the second demand, so the noise never fully cancels; that is exactly the jitter
conceded there, and Appendix~\ref{app:runavg} explains what the constant buys in
exchange. The recipe is due to Herbert Robbins and Sutton Monro, who published it in 1951.

\section{The running average}\label{app:runavg}

Appendix~\ref{app:rm} gave the general principle; here is the alternative, elementary
explanation, needing nothing beyond school algebra. Let $A_n$ be the average of the first
$n$ observations $x_1, \dots, x_n$, that is, $A_n = (x_1 + \cdots + x_n)/n$. Then
\[
n\,A_n \;=\; x_1 + \cdots + x_n \;=\; (n-1)\,A_{n-1} + x_n.
\]
In words: $n$ times the average is the total of all $n$ observations, and that total is
the total of the first $n-1$ of them, which is $(n-1)\,A_{n-1}$ (average times count),
plus the newest one. Dividing by $n$ and rearranging,
\[
A_n \;=\; A_{n-1} + \frac{1}{n}\,\bigl(x_n - A_{n-1}\bigr).
\]
In words: each new observation moves the average by $1/n$ of the surprise, the gap between
the newcomer and the average so far. A two-number check: the average of $10$ and $16$ is
$13$; a third observation of $22$ moves it to $13 + \tfrac{1}{3}\,(22 - 13) = 16$, which
is indeed the average of $10$, $16$, and $22$.

Now replace the shrinking $1/n$ by a small constant $\alpha$. The result is an average
that gently favors recent observations: each old observation's influence shrinks by a
factor $1 - \alpha$ with every new week. That is the right choice here, because the early
evidence was measured against a poor table $q$ and should fade. With that one
substitution, the running-average recipe \emph{is} the SARSA update~\eqref{eq:sarsa}: the
rule is nothing more exotic than a running average of the weekly evidence.

\end{document}
